\subsection{关于向量测度的向量值函数的积分}

命题 11. --- 设 F、G、H 为三个 Banach 空间， \(\Phi\) 为 F × G 到 H 中的连续双线性映射。设 m 为 T 上取值于 G 的可优超向量测度。则存在唯一的连续线性映射 I \(_{\Phi,\mathbf{m}}\) （由 \(\overline{\mathcal{L}}_{\mathrm{F}}^{1}(|\mathbf{m}|)\) 到 H 中），使得对每个 z ∈ F 与每个关于 \textbar m\textbar 可积的数值函数 h ，有 I \(_{\Phi,\mathbf{m}}(hz) = \Phi(z, \int h dm)\) 。此外，

\[
\left| \mathrm{I} _ {\Phi , \mathbf {m}} (\mathbf {f}) \right| \leqslant \| \Phi \| \int | \mathbf {f} | d | \mathbf {m} |\tag{4}
\]

对每个函数 \(\mathbf{f} \in \overline{\mathcal{L}}_{\mathrm{F}}^{1}(|\mathbf{m}|)\) 成立。

若存在这样的映射，它对 \(\mathbf{f}\) （ \(|\mathbf{m}|\) -可积集上的阶梯函数）所取的值就被唯一确定：因为已知这时可以写 \(\mathbf{f} = \sum_{i} \mathbf{a}_{i} \varphi_{\mathrm{X}_{i}}\) ，其中诸 \(\mathrm{X}_{i}\) 是 \(|\mathbf{m}|\) -可积且两两不相交的，而诸 \(\mathbf{a}_{i} \in \mathrm{F}\) （Ch. IV, §4, No.~9, 引理）。于是 \(\mathrm{I}_{\Phi, \mathbf{m}}(\mathbf{f})\) 的值就必须等于 \(\sum_{i} \Phi(\mathbf{a}_{i}, \int \varphi_{\mathrm{X}_{i}} d\mathbf{m})\) 。而今，我们有（No.~3, 命题 5）

\[
\left| \sum_ {i} \Phi (\mathbf {a} _ {i}, \int \varphi_ {\mathrm{X} _ {i}} d \mathbf {m}) \right| \leqslant \| \Phi \| \cdot \sum_ {i} | \mathbf {a} _ {i} | \cdot | \mathbf {m} | (\mathrm{X} _ {i}) = \| \Phi \| \int | \mathbf {f} | d | \mathbf {m} |,\tag{5}
\]

这首先表明 H 的元素 \(\sum_{i} \Phi(a_i, \int \varphi_{X_i} d\mathbf{m})\) 不依赖于 \(\mathbf{f}\) 写成 \(\sum_{i} a_i \varphi_{X_i}\) 的特殊形式，因此可以把它记作 \(I_{\Phi, \mathbf{m}}(\mathbf{f})\) 。立即可以验证，如此定义的映射 \(I_{\Phi, \mathbf{m}}\) 在空间 \(\mathcal{E}_F\) （由 \(|\mathbf{m}|\) -可积集上的阶梯函数组成的空间）上是线性的：因为，只需把两个函数 \(\mathbf{f}, \mathbf{g}\) （ \(\mathcal{E}_F\) 中的函数）写成 \(\mathbf{f} = \sum_{i} a_i \varphi_{X_i}\) 与 \(\mathbf{g} = \sum_{i} b_i \varphi_{X_i}\) 的形式，其中诸 \(|\mathbf{m}|\) -可积集 \(X_i\) 作成一个两两不相交的有限族（根据 Ch. IV, §4, No.~9 的引理，这是可能的）。于是不等式 (5) 表明 \(I_{\Phi, \mathbf{m}}\) 在 \(\mathcal{E}_F\) 上连续，而由于这个子空间在 \(\overline{\mathscr{L}}_F^1\) 中稠密（Ch. IV, §4, No.~10, 命题 19 的推论 1 以及 Ch. V, §1, No.~3），由此即得 \(I_{\Phi,m}\) 的存在性与唯一性，以及不等式 (4)。

我们把 \(I_{\Phi,\mathbf{m}}(\mathbf{f})\) 称为 f 关于 m 的积分（相对于双线性映射 \(\Phi\) ）；当双线性映射 \(\Phi\) 在点 \((\mathbf{x},\mathbf{y})\) 处的值记作 xy 时，我们将用 \(\int f dm\) 代替 \(I_{\Phi,\mathbf{m}}(\mathbf{f})\) 。

采用第 6 目的记号，积分 \(\int \langle \mathbf{f},\mathbf{g}\rangle d\mu\) 不是别的，正是 \(\mathrm{I}_{\Phi ,\mathbf{m}}(\mathbf{f})\) （取 \(\Phi (\mathbf{x},\mathbf{x}^{\prime}) = \langle \mathbf{x},\mathbf{x}^{\prime}\rangle\) 与 \(\mathbf{m} = \mathbf{g}\cdot \mu\) ）

推论. --- 若 m 与 \(m'\) 是 T 上取值于 G 的两个可优超测度，则 \(I_{\Phi,m+m'} = I_{\Phi,m} + I_{\Phi,m'}\) ，并且 \(I_{\Phi,\lambda m} = \lambda I_{\Phi,m}\) （对每个标量 \(\lambda\) ）。

第二个断言是直接的。第一个断言的意思是，对每个既 \(|m|\) -可积又 \(|m'|\) -可积的函数 f，

\[
\mathrm{I} _ {\Phi , \mathbf {m} + \mathbf {m} ^ {\prime}} (\mathbf {f}) = \mathrm{I} _ {\Phi , \mathbf {m}} (\mathbf {f}) + \mathrm{I} _ {\Phi , \mathbf {m} ^ {\prime}} (\mathbf {f}).\tag{6}
\]

函数 \(\mathbf{f}\) 是 \((|\mathbf{m}| + |\mathbf{m}'|)\) -可积的（Ch. V, §2, No.~2, 命题 3 的推论 1），从而更是 \((|\mathbf{m} + \mathbf{m}'|)\) -可积的，于是 (6) 的第一端确实有意义。为证明 (6)，只需对 \(\mathbf{f}\) 是 \((|\mathbf{m}| + |\mathbf{m}'|)\) -可积集上的阶梯函数的情形加以证明，因为这些函数所成之集在 \(\overline{\mathcal{L}}_{\mathrm{F}}^{1}(|\mathbf{m}| + |\mathbf{m}'|)\) 中稠密，且由 (4)，(6) 的两端在该空间中都连续。而对 \(\mathbf{f} = \mathbf{a}\varphi_{\mathrm{X}}\) （其中 \(\mathrm{X}\) 是 \((|\mathbf{m}| + |\mathbf{m}'|)\) -可积的），(6) 的两端都化为 \(\Phi(\mathbf{a}, \int \varphi_{\mathrm{X}} d\mathbf{m}) + \Phi(\mathbf{a}, \int \varphi_{\mathrm{X}} d\mathbf{m}')\) ，由此即得推论。

注记. --- 当 \(\mathbf{m}\) 具有形式 \(\mathbf{b}\mu\) （其中 \(\mathbf{b} \in G\) ， \(\mu\) 是 \(T\) 上的实测度）时，

\[
\mathrm{I} _ {\Phi , \mathbf {m}} (\mathbf {f}) = \int \Phi (\mathbf {f} (t), \mathbf {b}) d \mu (t)
\]

对每个函数 \(\mathbf{f} \in \mathcal{L}_{\mathrm{F}}^{1}(\mu)\) 成立，因为两端在该空间上连续，且当 \(\mathbf{f}\) 是 \(|\mu|\) -可积集上的阶梯函数时它们一致。

